%% ma: L12-B09-0006
%% lop: 12
%% chuong: 3
%% bai: 9
%% dang: 12.9.3
%% loai: TN4
%% muc: TH
%% dap_an: "A"
%% nguon: "(NBV)-ĐỀ SỐ 18-MỖI NGÀY 1 ĐỀ THI 2025.pdf (D18, MỖI NGÀY 1 ĐỀ - Đề số 18), Phần I Câu 12"
%% kien_thuc: "Bài 9 (khoảng tứ phân vị của mẫu số liệu ghép nhóm)"
%% trang_thai: chinh-thuc
%% ghi_chu: "Đáp án gốc A đúng (giải độc lập: n = 84, Q1 = 7,5 + 2/17 ≈ 7,618, Q3 = 8,5 + 5/26 ≈ 8,692, khoảng tứ phân vị 475/442 ≈ 1,0747 ≈ 1,1). Đề gốc ghi vào 10 năm học (thiếu chữ lớp) và không nêu cách làm tròn; bản kho sửa thành vào lớp 10 và thêm yêu cầu làm tròn đến hàng phần mười để chỉ có một đáp án. Trùng ý tưởng với dạng đã duyệt 12.9.3 (khoảng tứ phân vị), nhập chính thức."
\begin{ex}
Theo thống kê điểm trung bình môn Toán của một số học sinh đã trúng tuyển vào lớp $10$ năm học $2024$--$2025$ của Trường TH--THCS--THPT Lê Thánh Tông được kết quả như bảng sau:
\begin{center}
\begin{tabular}{|c|c|c|c|c|c|c|c|}\hline
Khoảng điểm & $[6{,}5;7)$ & $[7;7{,}5)$ & $[7{,}5;8)$ & $[8;8{,}5)$ & $[8{,}5;9)$ & $[9;9{,}5)$ & $[9{,}5;10)$\\\hline
Tần số & $7$ & $10$ & $17$ & $24$ & $13$ & $8$ & $5$\\\hline
\end{tabular}
\end{center}
Khoảng tứ phân vị của mẫu số liệu ghép nhóm trên (làm tròn đến hàng phần mười) là
\choice{\True $\Delta_Q=1{,}1$}{$\Delta_Q=1$}{$\Delta_Q=1{,}2$}{$\Delta_Q=0{,}6$}
\loigiai{Cỡ mẫu $n=7+10+17+24+13+8+5=84$.

Tứ phân vị thứ nhất ứng với $\dfrac n4=21$, thuộc nhóm $[7{,}5;8)$ (tần số tích lũy trước đó là $17$): $Q_1=7{,}5+\dfrac{21-17}{17}\cdot0{,}5=7{,}5+\dfrac2{17}\approx7{,}618$.

Tứ phân vị thứ ba ứng với $\dfrac{3n}4=63$, thuộc nhóm $[8{,}5;9)$ (tần số tích lũy trước đó là $58$): $Q_3=8{,}5+\dfrac{63-58}{13}\cdot0{,}5=8{,}5+\dfrac5{26}\approx8{,}692$.

Khoảng tứ phân vị $\Delta_Q=Q_3-Q_1\approx1{,}075\approx1{,}1$. Chọn A.}
\end{ex}
